% ==============================================================
\section{関連研究}\label{sec:related}
% ==============================================================

本研究は、頻出パターンのサポート変動を外部イベントに帰属させる問題を
空間次元に拡張する。この問題は複数の研究領域にまたがるため、
各領域の代表的手法とその限界を整理し、本研究の位置づけを明確にする。

\subsection{イベント帰属パイプライン}

頻出パターンのサポート変動を外部イベントに帰属させるフレームワーク~\cite{dsaa2025}は、
変化点検出・帰属スコアリング・置換検定・BH補正・重複排除の5ステップで構成される。
このパイプラインは\textbf{1次元のサポート時系列}のみを扱うため、
イベント効果が空間的に局在する場合に2つの限界がある：
(1) 全空間のサポートを集約することでイベント効果が希釈され検出力が低下する、
(2)「どの地域で効果があったか」という問いに答えられない。
\textbf{本研究}はこの1次元パイプラインを多次元サポート曲面に拡張し、
空間的近接度と空間構造を保存する置換検定を導入することで、
これらの限界を克服する。

\subsection{空間スキャン統計}

Kulldorff~\cite{kulldorff1997spatial}の空間スキャン統計は、
点過程データにおける空間クラスターを検出する代表的手法であり、
時空間拡張~\cite{kulldorff2001prospective}は円筒形の走査窓を使用する。
しかしこれらは\textbf{アイテムセットパターンを扱わず}、
検出された密集が\emph{なぜ}生じたかというイベントへの帰属には答えない。
本研究は、パターンのサポート\emph{変動}を検出するだけでなく、
その変動を既知の外部イベントに統計的に帰属させる点で根本的に異なる。

\subsection{密部分テンソルマイニング}

DenseAlert~\cite{shin2017densealert}やM-Zoom~\cite{shin2018mzoom}は
多次元テンソルにおける密ブロックをリアルタイムまたは事後的に検出する。
これらは密度\emph{異常}の検出に優れるが、
検出された密度変動と外部イベントの\textbf{統計的関連を検定する機能を持たない}。
すなわち「密である」ことは分かるが「なぜ密になったか」は不明のままである。
本研究は帰属スコアリングと置換検定を通じて、
変動の因果的候補（イベント）との関連を統計的に評価する。

\subsection{時空間パターンマイニング}

軌跡パターンマイニング~\cite{giannotti2007trajectory}や
空間コロケーションパターン~\cite{shekhar2001colocation}は、
空間的に共起するパターンの発見を対象とする。
これらはパターンの空間的分布を扱うが、
サポートの\textbf{時間的変動}やその原因帰属は対象としない。
本研究は時間と空間の両方でサポート変動を追跡し、
イベントとの帰属関係を検定する点で異なる。

\subsection{計量経済学のイベント・スタディ法}

MacKinlay~\cite{mackinlay1997event}のイベント・スタディ法は、
企業イベント（合併・増配等）の株価への影響を統計的に評価する確立された手法である。
しかし\textbf{単一のアウトカム変数}（異常収益率）を対象としており、
数千のパターン時系列を同時に検定する本研究の設定とはスケールが異なる。
また空間次元は考慮されない。

\subsection{サロゲートデータと置換検定}

Lancaster et al.~\cite{lancaster2018surrogate}は
時系列データに対するサロゲートデータ生成法を体系的に整理しており、
円形シフトは時間的自己相関を保存する代表的な手法である。
本研究の時空間置換検定は、この時間的サロゲート法を
\textbf{空間構造の保存に拡張}したものである。
具体的には、イベントの時刻を円形シフトしつつ空間スコープを固定することで、
空間自己相関構造を保ったまま帰無分布を生成する（\ref{sec:st_permtest}節で詳述）。

\subsection{本研究の位置づけ}

表~\ref{tab:related}に各関連手法との比較を示す。
本研究は、(1)~アイテムセットパターンのサポート変動を扱い、
(2)~空間次元を明示的にモデル化し、
(3)~統計的検定によるイベント帰属を行う、
唯一のフレームワークである。

\begin{table}[t]
\centering
\caption{関連手法との比較。
$\checkmark$: 対応、$\times$: 非対応。}
\label{tab:related}
\small
\begin{tabular}{lccc}
\toprule
手法 & パターン & 空間 & 帰属検定 \\
\midrule
1Dイベント帰属~\cite{dsaa2025}        & $\checkmark$ & $\times$     & $\checkmark$ \\
空間スキャン統計~\cite{kulldorff1997spatial} & $\times$     & $\checkmark$ & $\times$     \\
DenseAlert~\cite{shin2017densealert}  & $\times$     & $\checkmark$ & $\times$     \\
コロケーション~\cite{shekhar2001colocation} & $\checkmark$ & $\checkmark$ & $\times$     \\
イベント・スタディ~\cite{mackinlay1997event} & $\times$     & $\times$     & $\checkmark$ \\
\midrule
\textbf{本研究}                         & $\checkmark$ & $\checkmark$ & $\checkmark$ \\
\bottomrule
\end{tabular}
\end{table}
